% prezentace.tex -- Lekce 08: Funkce
% preamble-prez.tex -- sdílená preamble pro prezentace (beamer)
% Includuje se z ucitel/lekceXX/prezentace.tex jako:
%   % preamble-prez.tex -- sdílená preamble pro prezentace (beamer)
% Includuje se z ucitel/lekceXX/prezentace.tex jako:
%   % preamble-prez.tex -- sdílená preamble pro prezentace (beamer)
% Includuje se z ucitel/lekceXX/prezentace.tex jako:
%   \input{../spolecne/preamble-prez.tex}

\documentclass[aspectratio=169]{beamer}

% --- Jazyk a font ---
\usepackage{polyglossia}
\setmainlanguage{czech}
\usepackage{fontspec}

% --- Téma ---
\usetheme{default}
\usecolortheme{default}

\definecolor{microbit-blue}{HTML}{1E4D8C}
\definecolor{microbit-lightblue}{HTML}{D6E4F7}
\definecolor{code-bg}{HTML}{F0F0F0}

\setbeamercolor{frametitle}{bg=microbit-blue, fg=white}
\setbeamercolor{title}{fg=microbit-blue}
\setbeamercolor{block title}{bg=microbit-blue, fg=white}
\setbeamercolor{block body}{bg=microbit-lightblue}
\setbeamercolor{itemize item}{fg=microbit-blue}
\setbeamercolor{itemize subitem}{fg=microbit-blue!70}
\setbeamerfont{frametitle}{size=\large, series=\bfseries}
\setbeamertemplate{navigation symbols}{}
\setbeamertemplate{footline}{%
  \leavevmode%
  \hbox{%
    \begin{beamercolorbox}[wd=.5\paperwidth,ht=2.5ex,dp=1.125ex,leftskip=6pt]{author in head/foot}%
      \usebeamerfont{author in head/foot}\textcolor{gray}{\small Úvod do Pythonu na Micro:bitu}
    \end{beamercolorbox}%
    \begin{beamercolorbox}[wd=.5\paperwidth,ht=2.5ex,dp=1.125ex,rightskip=6pt,right]{date in head/foot}%
      \usebeamerfont{date in head/foot}\textcolor{gray}{\small\insertframenumber\,/\,\inserttotalframenumber}
    \end{beamercolorbox}}%
}

% --- Česky korektní uvozovky ---
\usepackage{csquotes}
\newcommand{\uv}[1]{\enquote{#1}}

% --- Zdrojový kód (Python) ---
\usepackage{listings}
\lstset{
  language=Python,
  basicstyle=\ttfamily\small,
  backgroundcolor=\color{code-bg},
  frame=single,
  framerule=0pt,
  xleftmargin=6pt,
  xrightmargin=6pt,
  breaklines=true,
  showstringspaces=false,
  keywordstyle=\color{microbit-blue}\bfseries,
  stringstyle=\color{teal},
  commentstyle=\color{gray}\itshape,
  literate=
    {á}{{\'a}}1 {é}{{\'e}}1 {í}{{\'i}}1 {ó}{{\'o}}1 {ú}{{\'u}}1
    {ů}{{\r{u}}}1 {č}{{\v{c}}}1 {š}{{\v{s}}}1 {ž}{{\v{z}}}1
    {Á}{{\'A}}1 {É}{{\'E}}1 {Í}{{\'I}}1 {Ó}{{\'O}}1 {Ú}{{\'U}}1
    {Č}{{\v{C}}}1 {Š}{{\v{S}}}1 {Ž}{{\v{Z}}}1 {ň}{{\v{n}}}1 {ř}{{\v{r}}}1,
}

% --- Zvýraznění pojmů ---
\newcommand{\pojem}[1]{\textbf{\textcolor{microbit-blue}{#1}}}
\newcommand{\pyinline}[1]{\texttt{#1}}

% --- Grafika a diagramy ---
\usepackage{tikz}

% --- Ostatní ---
\usepackage{graphicx}
\usepackage{hyperref}
\hypersetup{colorlinks=true, urlcolor=microbit-blue}


\documentclass[aspectratio=169]{beamer}

% --- Jazyk a font ---
\usepackage{polyglossia}
\setmainlanguage{czech}
\usepackage{fontspec}

% --- Téma ---
\usetheme{default}
\usecolortheme{default}

\definecolor{microbit-blue}{HTML}{1E4D8C}
\definecolor{microbit-lightblue}{HTML}{D6E4F7}
\definecolor{code-bg}{HTML}{F0F0F0}

\setbeamercolor{frametitle}{bg=microbit-blue, fg=white}
\setbeamercolor{title}{fg=microbit-blue}
\setbeamercolor{block title}{bg=microbit-blue, fg=white}
\setbeamercolor{block body}{bg=microbit-lightblue}
\setbeamercolor{itemize item}{fg=microbit-blue}
\setbeamercolor{itemize subitem}{fg=microbit-blue!70}
\setbeamerfont{frametitle}{size=\large, series=\bfseries}
\setbeamertemplate{navigation symbols}{}
\setbeamertemplate{footline}{%
  \leavevmode%
  \hbox{%
    \begin{beamercolorbox}[wd=.5\paperwidth,ht=2.5ex,dp=1.125ex,leftskip=6pt]{author in head/foot}%
      \usebeamerfont{author in head/foot}\textcolor{gray}{\small Úvod do Pythonu na Micro:bitu}
    \end{beamercolorbox}%
    \begin{beamercolorbox}[wd=.5\paperwidth,ht=2.5ex,dp=1.125ex,rightskip=6pt,right]{date in head/foot}%
      \usebeamerfont{date in head/foot}\textcolor{gray}{\small\insertframenumber\,/\,\inserttotalframenumber}
    \end{beamercolorbox}}%
}

% --- Česky korektní uvozovky ---
\usepackage{csquotes}
\newcommand{\uv}[1]{\enquote{#1}}

% --- Zdrojový kód (Python) ---
\usepackage{listings}
\lstset{
  language=Python,
  basicstyle=\ttfamily\small,
  backgroundcolor=\color{code-bg},
  frame=single,
  framerule=0pt,
  xleftmargin=6pt,
  xrightmargin=6pt,
  breaklines=true,
  showstringspaces=false,
  keywordstyle=\color{microbit-blue}\bfseries,
  stringstyle=\color{teal},
  commentstyle=\color{gray}\itshape,
  literate=
    {á}{{\'a}}1 {é}{{\'e}}1 {í}{{\'i}}1 {ó}{{\'o}}1 {ú}{{\'u}}1
    {ů}{{\r{u}}}1 {č}{{\v{c}}}1 {š}{{\v{s}}}1 {ž}{{\v{z}}}1
    {Á}{{\'A}}1 {É}{{\'E}}1 {Í}{{\'I}}1 {Ó}{{\'O}}1 {Ú}{{\'U}}1
    {Č}{{\v{C}}}1 {Š}{{\v{S}}}1 {Ž}{{\v{Z}}}1 {ň}{{\v{n}}}1 {ř}{{\v{r}}}1,
}

% --- Zvýraznění pojmů ---
\newcommand{\pojem}[1]{\textbf{\textcolor{microbit-blue}{#1}}}
\newcommand{\pyinline}[1]{\texttt{#1}}

% --- Grafika a diagramy ---
\usepackage{tikz}

% --- Ostatní ---
\usepackage{graphicx}
\usepackage{hyperref}
\hypersetup{colorlinks=true, urlcolor=microbit-blue}


\documentclass[aspectratio=169]{beamer}

% --- Jazyk a font ---
\usepackage{polyglossia}
\setmainlanguage{czech}
\usepackage{fontspec}

% --- Téma ---
\usetheme{default}
\usecolortheme{default}

\definecolor{microbit-blue}{HTML}{1E4D8C}
\definecolor{microbit-lightblue}{HTML}{D6E4F7}
\definecolor{code-bg}{HTML}{F0F0F0}

\setbeamercolor{frametitle}{bg=microbit-blue, fg=white}
\setbeamercolor{title}{fg=microbit-blue}
\setbeamercolor{block title}{bg=microbit-blue, fg=white}
\setbeamercolor{block body}{bg=microbit-lightblue}
\setbeamercolor{itemize item}{fg=microbit-blue}
\setbeamercolor{itemize subitem}{fg=microbit-blue!70}
\setbeamerfont{frametitle}{size=\large, series=\bfseries}
\setbeamertemplate{navigation symbols}{}
\setbeamertemplate{footline}{%
  \leavevmode%
  \hbox{%
    \begin{beamercolorbox}[wd=.5\paperwidth,ht=2.5ex,dp=1.125ex,leftskip=6pt]{author in head/foot}%
      \usebeamerfont{author in head/foot}\textcolor{gray}{\small Úvod do Pythonu na Micro:bitu}
    \end{beamercolorbox}%
    \begin{beamercolorbox}[wd=.5\paperwidth,ht=2.5ex,dp=1.125ex,rightskip=6pt,right]{date in head/foot}%
      \usebeamerfont{date in head/foot}\textcolor{gray}{\small\insertframenumber\,/\,\inserttotalframenumber}
    \end{beamercolorbox}}%
}

% --- Česky korektní uvozovky ---
\usepackage{csquotes}
\newcommand{\uv}[1]{\enquote{#1}}

% --- Zdrojový kód (Python) ---
\usepackage{listings}
\lstset{
  language=Python,
  basicstyle=\ttfamily\small,
  backgroundcolor=\color{code-bg},
  frame=single,
  framerule=0pt,
  xleftmargin=6pt,
  xrightmargin=6pt,
  breaklines=true,
  showstringspaces=false,
  keywordstyle=\color{microbit-blue}\bfseries,
  stringstyle=\color{teal},
  commentstyle=\color{gray}\itshape,
  literate=
    {á}{{\'a}}1 {é}{{\'e}}1 {í}{{\'i}}1 {ó}{{\'o}}1 {ú}{{\'u}}1
    {ů}{{\r{u}}}1 {č}{{\v{c}}}1 {š}{{\v{s}}}1 {ž}{{\v{z}}}1
    {Á}{{\'A}}1 {É}{{\'E}}1 {Í}{{\'I}}1 {Ó}{{\'O}}1 {Ú}{{\'U}}1
    {Č}{{\v{C}}}1 {Š}{{\v{S}}}1 {Ž}{{\v{Z}}}1 {ň}{{\v{n}}}1 {ř}{{\v{r}}}1,
}

% --- Zvýraznění pojmů ---
\newcommand{\pojem}[1]{\textbf{\textcolor{microbit-blue}{#1}}}
\newcommand{\pyinline}[1]{\texttt{#1}}

% --- Grafika a diagramy ---
\usepackage{tikz}

% --- Ostatní ---
\usepackage{graphicx}
\usepackage{hyperref}
\hypersetup{colorlinks=true, urlcolor=microbit-blue}


\title{Lekce 08: Funkce}
\subtitle{Úvod do Pythonu na Micro:bitu}
\date{}

\begin{document}

\begin{frame}
  \titlepage
\end{frame}

% ------------------------------------------------------------------
\begin{frame}[fragile]{Proč funkce?}

  Opakující se kód je problém:

  \bigskip
  \begin{lstlisting}
# bez funkce -- stejny kod 3x
if button_a.was_pressed():
    display.show(Image.HAPPY)
    sleep(500)
    display.clear()
if button_b.was_pressed():
    display.show(Image.HAPPY)
    sleep(500)
    display.clear()
  \end{lstlisting}

  \pause\bigskip
  \begin{block}{DRY — Don't Repeat Yourself}
    Opakující se kód dáme do \pojem{funkce}.
    Funkci pak voláme \textbf{kdekoliv, kolikrát chceme} — kód napíšeme jen jednou.
  \end{block}

  \bigskip
  Vzpomínáte na mikrovlnku z~lekce~1?\\
  Tlačítko „Start" = volání funkce. Nyní si tlačítka \textbf{vyrábíme sami}.
\end{frame}

% ------------------------------------------------------------------
\begin{frame}[fragile]{Definice funkce — \texttt{def}}

  \begin{lstlisting}
def blikni():
    display.show(Image.HAPPY)
    sleep(500)
    display.clear()
  \end{lstlisting}

  \bigskip
  \begin{itemize}
    \item \pyinline{def} — klíčové slovo „definuj funkci"
    \item \pyinline{blikni} — název funkce (pravidla jako pro proměnné)
    \item \pyinline{()} — závorky jsou povinné (i bez parametrů)
    \item \pyinline{:} a odsazení — stejně jako u \pyinline{if} nebo \pyinline{while}
  \end{itemize}

  \pause\bigskip
  \begin{alertblock}{Definice vs. volání}
    \pyinline{def blikni():} funkci \textbf{připraví}, ale \textbf{nespustí}.\\
    \pyinline{blikni()} funkci \textbf{spustí} — závorky jsou nutné i při volání!
  \end{alertblock}
\end{frame}

% ------------------------------------------------------------------
\begin{frame}[fragile]{Parametry a argumenty}

  \pojem{Parametr} = proměnná v~definici; \pojem{argument} = hodnota při volání:

  \bigskip
  \begin{lstlisting}
def pozdrav(jmeno):          # jmeno = parametr
    display.scroll("Ahoj, " + jmeno + "!")

pozdrav("Karel")             # "Karel" = argument
pozdrav("Jana")
  \end{lstlisting}

  \bigskip
  \begin{columns}[T]
    \begin{column}{0.46\textwidth}
      \begin{block}{Parametr}
        Lokální proměnná funkce.\\
        Existuje jen uvnitř.\\
        Definujeme ji v závorce.
      \end{block}
    \end{column}
    \begin{column}{0.46\textwidth}
      \begin{block}{Argument}
        Konkrétní hodnota při volání.\\
        Přiřadí se do parametru.\\
        Píšeme do závorek při volání.
      \end{block}
    \end{column}
  \end{columns}
\end{frame}

% ------------------------------------------------------------------
\begin{frame}[fragile]{\texttt{return} — návratová hodnota}

  Funkce může \textbf{vrátit} výsledek zpět volajícímu:

  \bigskip
  \begin{lstlisting}
import random

def hod_kostkou(pocet_sten):
    return random.randint(1, pocet_sten)

cislo = hod_kostkou(6)      # cislo = vysledek
display.show(str(cislo))
  \end{lstlisting}

  \bigskip
  \begin{itemize}
    \item \pyinline{return} ukončí funkci a vrátí hodnotu
    \item Kód za \pyinline{return} se \textbf{neprovede}
    \item Funkce bez \pyinline{return} vrátí \pyinline{None}
  \end{itemize}

  \pause\bigskip
  \begin{exampleblock}{Tip}
    Funkce s \pyinline{return} je jako kalkulačka — zadáš vstup, dostaneš výsledek.\\
    Funkce bez \pyinline{return} je jako tiskárna — provede akci, nic nevrátí.
  \end{exampleblock}
\end{frame}

% ------------------------------------------------------------------
\begin{frame}[fragile]{Refaktoring — přepis existujícího kódu}

  Teploměr z~lekce~06 — přepíšeme podmínky do funkce:

  \bigskip
  \begin{lstlisting}
from microbit import *

def zobraz_stav(teplota):
    if teplota > 27:
        display.show(Image.HAPPY)
    elif teplota < 18:
        display.show(Image.SAD)
    else:
        display.show(Image.MEH)

while True:
    zobraz_stav(temperature())
    sleep(2000)
  \end{lstlisting}

  \bigskip
  Hlavní smyčka je nyní \textbf{přehledná} — stačí jeden řádek.\\
  Funkci \pyinline{zobraz_stav} lze snadno znovu použít nebo upravit.
\end{frame}

% ------------------------------------------------------------------
\begin{frame}{Co jsme se naučili}
  \begin{itemize}
    \item[\checkmark] \pyinline{def nazev():} — definuje funkci
    \item[\checkmark] Funkce se spustí až při \textbf{volání} \pyinline{nazev()}
    \item[\checkmark] \pojem{Parametr} — proměnná v definici; \pojem{argument} — hodnota při volání
    \item[\checkmark] \pyinline{return hodnota} — funkce vrátí výsledek
    \item[\checkmark] Funkce bez \pyinline{return} vrací \pyinline{None}
    \item[\checkmark] \pojem{Refaktoring} — přepis kódu do funkcí pro přehlednost (DRY)
  \end{itemize}
\end{frame}

\end{document}
